%!TEX root = ../template.tex

\typeout{NT FILE chapter4.tex}%

\chapter{Language to Specify CRDTs}
\label{cha:language}

Chapters~2 and~3 established that automated and deductive verification address CRDT correctness from different angles, and
that neither is sufficient on its own. The framework at the centre of this dissertation reconciles the two not by
translating one tool's artifacts into the other's, but by introducing a dedicated language in which a CRDT is specified
once, and from which both a VeriFx implementation and a Why3 specification are derived automatically.

Five sections follow, covering both the language itself and the constraints the compiler enforces on any specification
written in it. The first, Section~\ref{sec:syntax}, covers the concrete syntax used to declare a CRDT, its state, its
functions, and how these are grouped into a module. Section~\ref{sec:language-interfaces} then turns to how the
\texttt{CvRDT} and \texttt{CmRDT} interfaces referenced by every module are themselves declared, together with the
algebraic axioms that can be associated with them. Sections~\ref{sec:language-contracts} and~\ref{sec:language-directives}
describe two complementary sets of restrictions: first, the conditions a module's declarations must satisfy to be accepted
as an implementation of its interface; and second, the directives, \texttt{[@vfx]}, \texttt{[@proof]}, and
\texttt{variant(…)}, that let a specification diverge, in controlled ways, between the VeriFx and Why3 backends. The
chapter closes, in Section~\ref{sec:language-types}, with a complete description of the type system introduced informally
in the sections before it.

\section{Syntax}
\label{sec:syntax}

A specification in the framework's language is built around two complementary constructs: interfaces,
which state the contract an implementation must satisfy, its types, functions, and the algebraic properties expected of
them, and modules, which provide a concrete implementation of one such contract for a specific CRDT. This separation is
deliberate as it lets the compiler validate a module purely by checking it against the shape declared by its interface,
Section~\ref{sec:language-contracts}, independently of how the module's logic is eventually compiled into VeriFx or WhyML.

A module begins with a header of the form:

\begin{verbatim}
module <Name> : <Interface>
  ...
end
\end{verbatim}

which names the CRDT being defined and states, via the \texttt{:} after the name, which interface it claims to implement,
for instance, \texttt{module GCounter\_CRDT : CvRDT}. This single line is what lets the backend decide, later in the
pipeline, whether the module should be compiled along the state-based (\texttt{CvRDT}) or operation-based (\texttt{CmRDT})
path, without the programmer having to state this intent anywhere else.

Inside a module, type declarations introduce the vocabulary used by the rest of the definition. A concrete type is
introduced with \texttt{type <name> = <type expression>}, most commonly used to define the module's own state
representation, \texttt{type payload = ...}. A type may also be declared without a body, \texttt{type <name>},
which introduces it as a generic parameter rather than a concrete definition. \texttt{GSet} declares \texttt{type elem}
before using it in \texttt{type payload = set elem}, leaving the element type unconstrained so that the resulting CRDT
can later be instantiated for any type. The type expressions available on the right-hand side of a \texttt{type}
declaration include the primitive types \texttt{integer} and \texttt{boolean}, the collection types \texttt{map<K, V>}
and \texttt{set T}, tuples \texttt{(T1, T2)}, records \texttt{\{ field1: T1, field2: T2, ... \}}, and variants, written
either as a plain enumeration \texttt{A | B | C} or, when a case needs to carry data, as \texttt{A: T1 | B: T2},
Section~\ref{sec:language-types}.

Function declarations follow the keyword \texttt{val}, and share a single general shape:

\begin{verbatim}
val <name> (<parameters>) : <return type> = <body>
\end{verbatim}

Parameters sharing a type can be grouped under a single annotation, as in \texttt{(a b: payload)}.
The same pattern extends to any number of parameters, which keeps functions whose parameters repeat a type,
\texttt{merge}, \texttt{compare}, and \texttt{equals}, among others, concise to declare. When a function takes no
parameters, as is the case for the constructor \texttt{create}, the parentheses are omitted entirely, and the declaration
reduces to \texttt{val create : <return type> = <body>}. A function body may optionally be preceded by a
\texttt{requires \{ <condition> \}} clause, placed between the return type and the body itself. This expresses a
precondition on the function's inputs, Section~\ref{sec:language-directives}.

To make this concrete, Listing~\ref{lst:gcounter} gives the complete definition of \texttt{GCounter\_CRDT}, the running
example for the rest of this chapter.

\begin{lstlisting}[caption={GCounter\_CRDT module}, label={lst:gcounter}]
module GCounter_CRDT : CvRDT
  type payload = integer

  val create : payload = { payload : 0 }

  val merge (a b: payload) : payload =
    max(a, b)

  val compare (a b: payload) : boolean =
    a == b

  val equals (a b: payload) : boolean =
    compare(a, b) && compare(b, a)

  val increment (v: integer, a: payload) : payload =
    requires { v > 0 }
    a + v
end
\end{lstlisting}

One detail in this listing is easy to miss but worth explaining, since nothing in the syntax itself points to it.
Although \texttt{payload} is declared as a plain \texttt{integer}, the body of \texttt{create} is written as a record
literal, \texttt{\{ payload : 0 \}}, rather than simply \texttt{0}. This is not an inconsistency in the language but a
direct consequence of how every module's runtime state is represented internally. Regardless of whether a CRDT's
\texttt{payload} is itself a record, as it is for \texttt{PNCounter\_CRDT} or \texttt{TwoPSet}, whose \texttt{payload}
already has named fields, the framework always wraps a module's state in an implicit record with a \texttt{payload} field,
mirroring the two-module structure (\texttt{type t = \{ payload: ... \}}) that the Why3 backend generates for every CRDT
(Chapter~5). Writing \texttt{create} as a record literal is therefore not an arbitrary convention but a syntactic
consequence of this internal representation. It is the only place in a module where the shape of the compiler's internal
state representation surfaces in the code the programmer writes. Every other function operates on \texttt{payload}
directly, with no need to mention this wrapper.

\section{The CvRDT and CmRDT Interfaces}
\label{sec:language-interfaces}

Every module in the language is written to satisfy one of two interfaces, \texttt{CvRDT} for state-based CRDTs and
\texttt{CmRDT} for operation-based ones. Both are regular interfaces in the sense of Section~\ref{sec:syntax}. Nothing in the
compiler treats the identifiers \texttt{CvRDT} or \texttt{CmRDT} specially during parsing. What is specific to these two,
besides being the only two included in the framework, is that the compiler checks any interface carrying either of these
exact names for a fixed minimum signature, and that the backend inspects which of the two a module implements to decide
how to compile it, Section~\ref{sec:language-contracts}.

\subsection{Interface Declarations}
\label{sub:interface-decl}

An interface is declared with:

\begin{verbatim}
interface <Name>[@proof]
  ...
end
\end{verbatim}

The \texttt{[@proof]} tag is optional and related only to VeriFx, and its presence decides whether the generated proof
trait for this interface includes the base convergence proof, the structural \texttt{is\_a\_CvRDT}/\texttt{is\_a\_CmRDT}
obligation establishing the module as a genuine join-semilattice, or, for \texttt{CmRDT}, that its operations commute.
This is the same proof defined in VeriFx's own practical CRDT library. The framework's generated
\texttt{CvRDTProof}/\texttt{CmRDTProof} traits contain the same proof, so a module proven correct through this pipeline
carries the same guarantees as one written directly against VeriFx's own traits. Between \texttt{interface} and
\texttt{end}, only two kinds of declarations are allowed: types, written \texttt{type <name>} with no body, and function
signatures, written \texttt{val <name> (<parameters>) : <return type>} with no body either. An interface states what a
conforming module must provide. Listing~\ref{lst:interfaces} gives the \texttt{CvRDT} and \texttt{CmRDT} interfaces
exactly as declared in the framework:

\begin{lstlisting}[caption={The \texttt{CvRDT} and \texttt{CmRDT} interfaces}, label={lst:interfaces}]
interface CvRDT[@proof]
  type payload

  val create () : payload
  val merge (a b: payload) : payload
  val compare (a b: payload) : boolean
  val equals (a b: payload) : boolean

  axiom commutative(merge)
  axiom idempotent(merge)
  axiom associative(merge)
  axiom equivalence(compare)
end

interface CmRDT[@proof]
  type payload
  type operation

  val create () : payload
  val execute (op: operation, state: payload) : payload
  val compare (a b: payload) : boolean
  val equals (a b: payload) : boolean

  axiom op_commutative(execute)
end
\end{lstlisting}

The two interfaces share almost the same shape. Both declare a \texttt{payload} type, a \texttt{create} with no parameters,
and \texttt{compare}/\texttt{equals}. They diverge in exactly the operation each was designed around. \texttt{CvRDT}
exposes \texttt{merge}, taking two states and returning their combination, the same join operation described in
Section~\ref{sub:state_convergence}, now expressed as the signature every \texttt{CvRDT} module has to provide.
\texttt{CmRDT} instead declares \texttt{type operation} and a single \texttt{execute}, taking an operation and the current
state and returning the state after that operation has been applied. VeriFx presents this behavior splitting it into a
\texttt{prepare} step that turns an operation into a message and an \texttt{effect} step that applies a received message
to the state, Section~\ref{sub:modeling}. The language expresses this as a single function instead. A specification simply
states the result of executing an operation on a state, and it is the VeriFx backend's job to derive the
\texttt{prepare}/\texttt{effect} split from it, (Chapter~5). Why3, in turn, uses \texttt{execute} directly, with no
equivalent split. The result is that a single \texttt{execute} definition is enough to target both backends, at the cost
of that definition needing to be split apart again later, invisibly to the programmer, whenever the VeriFx backend needs
it.

\subsection{Axioms}
\label{sub:axioms}

Alongside its types and functions, an interface can declare axioms, properties that every module implementing the
interface is expected to satisfy. An axiom can be declared in one of two ways. The first names one of five built-in
properties:

\begin{verbatim}
axiom <kind>(<function>)
\end{verbatim}

where \texttt{<function>} is the function the property applies to. The compiler turns each of these directly into
automated proof obligations for VeriFx (Chapter~5), which is why this particular form is restricted to a fixed vocabulary
rather than left open: it recognises exactly the five properties below, and rejects any other name:

\begin{itemize}
  \item \texttt{commutative(merge)}: merging \texttt{x} with \texttt{y} produces the same result as merging \texttt{y}
  with \texttt{x}.
  \item \texttt{idempotent(merge)}: merging a state with itself leaves it unchanged.
  \item \texttt{associative(merge)}: merging three states gives the same result regardless of how the merges are grouped.
  \item \texttt{equivalence(compare)}: the \texttt{equals} method derived from \texttt{compare}, Section~\ref{sec:syntax}, agrees with
  the notion of equality between two states. Without this, \texttt{compare} could describe some relation weaker than
  genuine equivalence, and \texttt{equals} would no longer identify the same states \texttt{==} does.
  \item \texttt{op\_commutative(execute)}: applying two operations to the same state in either order leaves the state in
  an equal, reachable result, the operation-based counterpart of \texttt{commutative} for \texttt{merge}.
\end{itemize}

Each of these compiles, for VeriFx, into a fixed proof template built around the named function. For instance,
\texttt{commutative(merge)} always produces a proof stating that \texttt{x.merge(y)} and \texttt{y.merge(x)} are equal
and reachable for any two reachable, compatible \texttt{x} and \texttt{y}. This is why the set of properties is closed
rather than extensible. Each one corresponds to a hand-written proof shape in the backend, and allowing an arbitrary
property name would mean admitting a proof the backend has no template for.

The second form lets the programmer state a property the five built-in properties do not cover, by writing the condition
directly:

\begin{verbatim}
axiom <name>: <boolean expression>
\end{verbatim}

which attaches an arbitrary boolean formula, written using the same expression syntax as function bodies, directly to the
interface under a name of the programmer's choosing. This form exists in the grammar and is fully supported by the type
checker, but none of the ten CRDTs implemented so far has needed it. The five built-in properties have so far been
sufficient to state every property these specifications require.

The five built-in properties can also appear without this closing on an interface at all, directly attached to a function
declared inside a module rather than an interface, when that function has no body of its own:

\begin{verbatim}
val <n> (<parameters>) : <return type>
  axiom <kind>
  axiom <kind>
  ...
\end{verbatim}

\texttt{NestedMap}, Appendix~\ref{app:crdt_specifications}, Listing~\ref{lst:app_nestedmap}, uses exactly this for
\texttt{merge\_values} and \texttt{compare\_values}, two functions it declares but never defines, only assumes satisfy
\texttt{commutative}, \texttt{idempotent}, and \texttt{associative} for the former, and \texttt{equivalence} for the
latter. This is what lets \texttt{NestedMap} stay generic over any value type: rather than committing to one concrete way
of merging values, it states the properties a merge for that value type would need to have, and leaves the definition
itself to whichever value type the module is eventually used with. The function name is omitted here, unlike in an
interface's \texttt{axiom <kind>(<function>)}, because only one function is in scope for it to describe. Internally, both
compile to the exact same representation, they differ only in how they are written.

\section{Meeting an Interface's Requirements}
\label{sec:language-contracts}

A module is only accepted by the compiler once it has been checked against the interface named in its header. This
checking happens in two layers. The first applies to any interface, and simply confirms that a module provides everything
the interface declares. The second is stricter, but only applies to the two interfaces the framework actually ships with.
Naming an interface \texttt{CvRDT} or \texttt{CmRDT} carries an additional obligation of its own.

\subsection{Validating Modules}
\label{sub:conformance}

For every type the interface declares, the module must declare a type of the same name. For every function signature the
interface declares, the module must provide a function of the same name, with the same number of parameters and a
compatible return type. These checks are independent of one another and fail the moment they are not met, each with a
message naming exactly what is missing or mismatched: \texttt{Module '\%s' missing type '\%s' present in interface '\%s'.}
for a missing type, \texttt{Module '\%s' missing function '\%s'.} when a function is absent altogether,
\texttt{Function '\%s' has wrong number of arguments.} for an arity mismatch, and \texttt{Function '\%s' return type does
not respect the interface's.} when the return type disagrees.

The second layer applies only to interfaces named exactly \texttt{CvRDT} or \texttt{CmRDT}, and is checked once, when the
interface itself is declared, rather than separately for every module that later implements it. An interface named
\texttt{CvRDT} must declare at least \texttt{create}, \texttt{merge}, \texttt{compare}, and \texttt{equals} among its
functions. One named \texttt{CmRDT} must declare at least \texttt{create}, \texttt{execute}, \texttt{compare}, and
\texttt{equals}. An interface with either name that omits one of these is rejected outright, with
\texttt{Interface '\%s' must declare at least \%s.} naming what was expected.

A custom axiom, Section~\ref{sec_langauge-interfaces}, declared on an interface is also carried over to every module that implements it. When a
module is checked, each \texttt{axiom <name>: <expression>} on its interface is type-checked a second time, this time
against the module's own concrete \texttt{payload} type, and the resulting formula is added to the module as an assumption.
A module therefore never needs to restate a custom axiom itself. Satisfying the interface is enough for the axiom to
apply to it.

\subsection{State Representation and the \texttt{payload}/\texttt{t} Distinction}
\label{sub:payload-t}

The return-type check described above allows for one exception. A module may declare its own \texttt{type t}, in addition
to \texttt{payload}, to represent its state internally, as \texttt{VGCounter} and \texttt{UWMap} do. When it does, its
functions are allowed to return \texttt{t} in place of \texttt{payload} and still satisfy the interface. Without this
allowance, a module that needed a richer internal representation than \texttt{payload} alone, a vector indexed by replica
for \texttt{VGCounter}, or a pair of maps tracking added and removed elements for \texttt{UWMap}, would have no way to
express that representation without also changing what \texttt{payload} itself means to the outside world, and every other
module composing it would have to change along with it.

This same mechanism also explains how one CRDT can be built out of another. Whenever a module declares a function
literally named \texttt{create}, whether it uses \texttt{type t} or works with \texttt{payload} directly, the compiler
automatically registers two functions other modules can call on it from outside, \texttt{<Module>.create} and
\texttt{<Module>.get\_payload}. \texttt{PNCounter\_CRDT} relies on exactly this when it is built out of two
\texttt{GCounter\_CRDT} instances. Its own \texttt{create} calls \texttt{GCounter\_CRDT.create} twice, once for each of
its two counters, and its \texttt{merge} reads their values back out through \texttt{GCounter\_CRDT.get\_payload}.
\texttt{PNCounter\_CRDT} never needs to know that \texttt{GCounter\_CRDT}'s \texttt{payload} happens to be a plain integer
rather than something more elaborate. It only ever sees \texttt{GCounter\_CRDT} through these two functions. This is, in
the end, why composing CRDTs in the language works at all. Every module with a \texttt{create} carries this pair of
functions along with it automatically, without the programmer composing two CRDTs ever having to ask for them.

\section{Backend-Specific Directives}
\label{sec:language-directives}

Most of a specification compiles the same way for both backends. Neither backend needs to know the other one exists. A few
things do not work this way though, and the language has four directives for them. \texttt{[@proof]}, on an interface,
was already covered in, Section~\ref{sec:language-interfaces}, since it belongs to the interface as a whole rather than to one declaration inside it.
This section covers the other three: \texttt{[@vfx]} on a function, \texttt{[@vfx class: ...]} on a type, and the
\texttt{[@vfx]} version of \texttt{requires}. It also covers \texttt{variant(...)}, which is not tied to VeriFx in the
same way, but fits here for a reason explained in Section~\ref{sub:variant}.

\subsection{VeriFx Functions}
\label{sub:vfx-functions}

A function declaration can carry a simple \texttt{[@vfx]} tag, placed right after its parameters:

\begin{verbatim}
val <name> (<parameters>) [@vfx] : <return type> = <body>
\end{verbatim}

A function tagged this way only exists for VeriFx, and Why3 does not see it, as if it had never been declared. This is
needed because not every function a CRDT uses is something Why3's proof needs, some functions exist solely to support
VeriFx's own proof machinery. \texttt{OpTwoPSet}, for example, declares \texttt{preRemove}, \texttt{enabledSrc},
\texttt{forall\_op}, and \texttt{filter} this way. These help VeriFx prove convergence for a set-based CmRDT, while
remaining unused on the Why3 side of the same specification. Five names are treated specially when tagged like this:
\texttt{enabledSrc}, \texttt{enabledDown}, \texttt{compatibleS}, \texttt{reachable}, and \texttt{compatible}. These are
the same methods \texttt{CvRDT} and \texttt{CmRDT} already give a default implementation for, Section~\ref{sub:modeling}, so a module
that supplies one of them is overriding that default, and the generated VeriFx code writes \texttt{override} instead of
a plain \texttt{def}.

\subsection{Shaping the VeriFx Representation}
\label{sub:vfx-class}

The same \texttt{[@vfx ...]} syntax can also attach to a type declaration, where it can ask the VeriFx backend for one of
three different things, depending on what follows \texttt{[@vfx}.

The first extracts the type into its own standalone class:

\begin{verbatim}
type <n> = { ... } [@vfx class: <fn1>, <fn2>, ...]
\end{verbatim}

\texttt{RGA}, Appendix~\ref{app:crdt_specifications}, Listing~\ref{lst:app_rga}, needs a Lamport clock to order concurrent
insertions, so it declares \texttt{type lamport\_clock}, along with functions like \texttt{smaller\_clock},
\texttt{increment\_clock}, and \texttt{sync}, each marked \texttt{[@vfx]} as in the previous subsection. On their own,
these would just become VeriFx-only functions on the \texttt{RGA} class. The \texttt{[@vfx class: ...]} directive on
\texttt{lamport\_clock} changes that: it pulls the five named functions out into their own \texttt{LamportClock} class.
This way the generated VeriFx code models a Lamport clock the way someone writing VeriFx directly would, as its own type
with its own methods, rather than as several functions defined directly on \texttt{RGA}. What the compiler requires is
the \texttt{class:} part, so \texttt{[@class: ...]} alone works exactly the same.

The second overrides the VeriFx type the backend would otherwise derive for a field, replacing it with whatever type
expression follows \texttt{[@vfx}:

\begin{verbatim}
type t = { <field>: <type expression>, ... } [@vfx <VeriFx type>]
\end{verbatim}

\texttt{VGCounter}, Appendix~\ref{app:crdt_specifications}, Listing~\ref{lst:app_vgcounter}, represents its state as
\texttt{type t = \{ array: map<integer, integer> \}}, a map from replica number to that replica's own count. Translated
automatically, this would become a generic VeriFx \texttt{Map[Int, Int]}. Tagging the declaration
\texttt{[@vfx Vector[Int]]} instead tells the backend to represent \texttt{array} as a \texttt{Vector[Int]}, matching how
a vector-based counter is normally modelled directly in VeriFx, and giving access to vector-specific operations, such as
\texttt{.zip} and \texttt{.map}, used later when compiling \texttt{merge} and \texttt{compare} for this CRDT (Chapter~5).

The third reuses an existing CRDT module's own VeriFx class for a field, instead of deriving a new one:

\begin{verbatim}
type payload = { ..., <field>: <type expression>, ... } [@vfx <Module>: <field>]
\end{verbatim}

\texttt{RGA}'s \texttt{payload} carries an \texttt{added} and a \texttt{removed} set of elements, the exact shape already
defined by \texttt{OpTwoPSet}, Appendix~\ref{app:crdt_specifications}, Listing~\ref{lst:app_optwopset}. Rather than
deriving a second, unrelated representation for the same pair of sets, \texttt{[@vfx OpTwoPSet: vertices]} tells the
backend to represent them as an instance of \texttt{OpTwoPSet}'s own generated VeriFx class, exposed on \texttt{RGA}
under the field name \texttt{vertices}.

\subsection{VeriFx Preconditions}
\label{sub:vfx-requires}

A \texttt{requires} clause, Section~\ref{sec:syntax}, can also be marked \texttt{[@vfx]}:

\begin{verbatim}
requires { <condition> } [@vfx]
\end{verbatim}

On the VeriFx side this behaves exactly like an ordinary \texttt{requires}, becoming the same \texttt{if}/\texttt{else}
guard around the function body. On the Why3 side it disappears completely, no contract is generated at all, unlike an
ordinary \texttt{requires}, which becomes a real Why3 precondition. \texttt{RGA} does this for \texttt{successor}, whose
precondition, \texttt{lookup(e, a)}, is only declared \texttt{[@vfx]}. This lets VeriFx narrow down the states its solver
explores, without needing Why3 to prove a precondition that is not actually necessary for the rest of the specification.

\subsection{Proving Termination with \texttt{variant}}
\label{sub:variant}

A recursive function can be given a termination measure with \texttt{variant(...)}, placed after its return type:

\begin{verbatim}
val <name> (<parameters>) : <return type>
  variant (<expression>, ...)
  = <body>
\end{verbatim}

Why3 needs this to prove that a recursive function actually terminates: the expression given has to strictly decrease,
in a way that cannot decrease forever, on every recursive call. \texttt{VGCounter}'s value computation uses
\texttt{variant (size - replica)} for this. VeriFx has nothing like this requirement, a \texttt{variant} clause simply
does not affect what gets generated for it.

For a function like \texttt{VGCounter}'s, this is enough, since the decreasing quantity is already a parameter with a
clear meaning. \texttt{RGA}'s \texttt{find\_insert\_position} needs more than that. Nothing in its own parameters already
decreases as it walks along the structure, so the specification adds one just for this purpose, an explicit \texttt{fuel}
parameter, decremented by one on every recursive call, with \texttt{variant (fuel)}:

\begin{verbatim}
val find_insert_position (stamp: lamport_clock, e left right: elem,
                           a: payload, fuel: integer) : payload
  variant (fuel)
  =
    if fuel <= 0 then
      ...
    else if smaller_clock(stamp, right.timestamp) then
      find_insert_position(stamp, e, right, successor(right, a), a, fuel - 1)
    else
      ...
\end{verbatim}

Why3 needs the \texttt{if fuel <= 0 then ...} branch as a fallback for when \texttt{fuel} reaches zero, even though it
is not meant to run in practice. \texttt{fuel} itself has no meaning of its own beyond giving Why3 something it can prove
decreases. The backend recognises this exact pattern, a parameter that only counts down inside an \texttt{if <param> <= 0
then ... else if ... then recurse(..., <param> - 1) else ...} shape, and removes it completely when compiling for VeriFx.
The \texttt{fuel} parameter is removed, along with the base case, and only the \texttt{else if}/\texttt{else} branches are
left, as an ordinary recursive function:

\begin{lstlisting}[caption={\texttt{find\_insert\_position}, as generated for VeriFx}, label={lst:find-insert}]
private def find_insert_position(stamp: LamportClock, e: Elem[V], left: Elem[V], right: Elem[V]): RGA[V] = {
  if (stamp.smallerClock(right.timestamp)) {
    this.find_insert_position(stamp, e, right, this.successor(right))
  } else {
    val newClock = if (this.clock.smallerClock(e.timestamp)) {
                     e.timestamp
                   } else {
                     this.clock
                   }
    new RGA(newClock, this.vertices.add(e), this.edges.add(e, right).add(left, e), this.bottom_value)
  }
}
\end{lstlisting}

Without this, \texttt{fuel} would have to be present on the VeriFx side too: an argument with no purpose there, one the
programmer calling \texttt{find\_insert\_position} would have to supply some value for on every call, just to satisfy a
termination requirement that belongs only to Why3.

\section{Type System}
\label{sec:language-types}

Section~\ref{sec:syntax} introduced the type expressions available in the language, integer, boolean, set, map, tuples, records, and
variants, without describing how the compiler checks them. This section addresses that gap, how a type written by the
programmer is resolved into an internal representation the compiler can compare and reason about, how expressions are
checked against that representation, and how the language determines a type in the one case where it cannot be read
directly from the syntax, calls to \texttt{set.*} and \texttt{map.*}.

\subsection{Types}
\label{sub:types}

Internally, every type expression the programmer writes is resolved into one of a small, fixed set of internal
representations. \texttt{integer} and \texttt{boolean} become two fixed built-in types. A named type that is not one of
these two is resolved differently depending on what it refers to, a type declared without a body, such as \texttt{GSet}'s
\texttt{type elem}, becomes an abstract type, possessing no structure beyond its own name, while a type declared with a
body, such as \texttt{type payload = ...}, is resolved to the structure that appears on its right-hand side, and the name
itself is retained as an alias pointing to that structure. \texttt{map<K, V>} and \texttt{set T} become structured types
parameterised by the types of their keys, values, or elements, themselves resolved the same way, recursively. Tuples,
records, and variants follow the same pattern, \texttt{(T1, T2)} becomes a pair of resolved types,
\texttt{\{ field1: T1, ... \}} becomes a list of field names paired with their resolved types, and a variant becomes
either a plain list of constructor names, for \texttt{A | B | C}, or a list of constructor names each paired with the
type of data it carries, for \texttt{A: T1 | B: T2}.

Because a name as \texttt{payload} is only an alias, comparing two types for compatibility cannot stop at the name, two
types can only be meaningfully compared once every alias they contain has been resolved to its underlying structure. The
compiler performs this resolution on demand, wherever two types need to be compared, rather than once up front, an alias
is expanded precisely when a given check requires looking past it. This is also why an abstract type such as \texttt{elem}
behaves differently from every other type, it has no structure to expand into, and is therefore equal only to itself.

\subsection{Typing Expressions}
\label{sub:typing-expressions}

Every expression within a function body is type-checked against the types visible at that point in the program, the
parameters in scope, the functions and records declared so far, and, where applicable, the type expected of the expression
by its surrounding context. Checking an expression produces both its type and an internal, fully-typed representation of
it, the same representation the backend later compiles from, leaving no aspect of an expression's meaning implicit beyond
this point.

Most of the rules are what a reader familiar with a typed language would expect. Arithmetic operators require both sides
to be \texttt{integer} and produce an \texttt{integer}. The ordering comparisons, \texttt{<} and \texttt{<=} among them,
impose the same requirement. Equality, \texttt{==}, requires only that both sides resolve to the same type once aliases
are expanded. The boolean connectives require both sides to be \texttt{boolean}. An \texttt{if} requires its condition to
be \texttt{boolean} and its two branches to agree on a single type, which becomes the type of the \texttt{if} expression
as a whole. A function call, whether to a value declared within the same module or to one exposed by another module
through the \texttt{create}/\texttt{get\_payload} mechanism described in Section~\ref{sec:language-contracts}, is checked for both the correct
number of arguments and the correct type of each.

A \texttt{match} expression corresponds to what other functional languages term pattern matching, and can match on more
than one expression at once, as \texttt{RGA}'s \texttt{op\_commutative} lemma does when it matches on a pair of operations
together, \texttt{op1} and \texttt{op2}, to establish commutativity case by case. Each branch's pattern is checked
against the structure of the variant type it matches, a constructor with no arguments, such as \texttt{Inc}, for a variant
declared as a simple enumeration, or a constructor applied to one variable for each piece of data it carries, such as
\texttt{Add v}, for a variant declared with \texttt{A: T1 | B: T2}. All branches are required to return the same type,
determined from the first branch checked.

Constructing a record is the one construct whose type is not stated anywhere within the expression itself. The literal
\texttt{\{ field1: e1, field2: e2 \}} names no type, only fields. The compiler resolves this by searching every record
type declared throughout the specification for the one whose field names match exactly, and using that record's field
types to check each of the given field expressions in turn. This matching is based purely on field names, not on the types
those fields hold, so two records that happen to declare the exact same field names cannot be told apart by it. This is
not a hypothetical case, \texttt{TwoPSet}, \texttt{OpTwoPSet}, and \texttt{UWMap}'s own internal \texttt{type t} all
declare a record with exactly the fields \texttt{added} and \texttt{removed}, differing only in what those fields actually
contain. Because nothing in the specification's correctness is meant to depend on the order in which a hash table happens
to be traversed, this is best understood as a possible source of fragility in the mechanism, rather than a guaranteed
property, one that has not surfaced as a problem for the ten CRDTs implemented so far, but that a more principled version
of this check, comparing field types as well as field names, would remove entirely. The sole exception to name-based
matching is the single-field literal \texttt{\{ payload : e \}}, used by \texttt{create} whenever a module's payload is
not itself a record, as discussed in Section~\ref{sec:syntax}.

\subsection{Resolving Types for Collection Calls}
\label{sub:collections}

\texttt{set} and \texttt{map} operations are shared across every CRDT in the specification, but the concrete element,
key, and value types are different in each module, an abstract element type in one CRDT, and a pair of tagged values in
another. A call such as \texttt{set.add(v, s)} or \texttt{map.get(k, m)} must therefore be typed according to where it
occurs, and the compiler resolves this for each call individually, rather than assuming one type shared by every use of
\texttt{set} or \texttt{map} in the specification.

In cases where an argument already carries the type required, the compiler reads it directly from that argument.
\texttt{set.add(v, s)} infers the element type from \texttt{v} itself, and then verifies that \texttt{s} is a set of that
same element type. \texttt{map.get(k, m)} infers the key and value types from \texttt{m}, and then checks \texttt{k}
against the key type thereby obtained. Some calls have no such argument to read from, \texttt{set.empty()} takes no
arguments, and \texttt{map.const(default)} supplies only a value, never a key. For these cases, the compiler falls back
on the type expected of the whole expression by its surrounding context, the same \texttt{expected} type threaded
throughout the rest of the type checker, for instance the declared type of the record field the call provides the value
for. If no expected type is available at that point, the compiler reports an error rather than attempting to guess.

\texttt{map.set(k, v, m)} illustrates why this inference is handled per operation rather than through a single, uniform
rule. It first attempts to read the key and value types from \texttt{m}, as \texttt{map.get} does, and only falls back
to inferring them from \texttt{k} and \texttt{v}, checking \texttt{m} against the resulting map type, if that first
attempt fails. Either direction is sufficient to type the call correctly, and attempting the more informative one first
avoids depending on context that is not always available. This per-operation handling is what lets \texttt{set} and
\texttt{map} behave as a single, shared vocabulary across every CRDT in the specification, rather than requiring a
distinct set of operations for every element, key, or value type a given module requires.
